% =====================================================================
\section{Introduction}
\label{sec:intro}
% =====================================================================

The negatively charged nitrogen-vacancy (N-V) centre in diamond is a solid-state electron spin that can be initialised
and read out optically at room temperature, which makes it a magnetic field sensor with nanoscale spatial resolution
\cite{balasubramanian2008,maze2008,doherty2013nv}. In optically detected magnetic resonance (ODMR) magnetometry the
field to be measured shifts the spin resonance, and it is read from the position of that resonance, found by
recording the photoluminescence at a sequence of microwave drive frequencies within a search window.

Choosing these drive frequencies well is an active field. Non-adaptive undersampling measures only a subset of
the window and reconstructs the line from it, either by fitting a Lorentzian to randomly chosen drive frequencies
\cite{luo2023} or by compressed sensing in a sparse spectral basis \cite{haim2025}. Bayesian experiment design
instead keeps a posterior over the parameters, here as a sequential Monte Carlo (SMC) particle cloud, and chooses
each next measurement to maximise its expected utility, such as the information it is expected to gain
\cite{granade2012}. In N-V Ramsey magnetometry, adaptive Bayesian schemes set the readout phase or the precession
time from the outcomes so far
\cite{bonato2016,santagati2019,mcmichael2021ramsey}; in ODMR, the same approach chooses the next drive frequency and
reaches a given precision more than an order of magnitude faster than a sweep
\cite{dushenko2020sbed,mcmichael2021optbayesexpt}.

In recent years, learning the design rule itself has become a particularly active direction. Early work
learned policies for adaptive phase estimation \cite{hentschel2010,lumino2018}. Reinforcement learning
has since been used to find the controls that make a quantum sensor most sensitive \cite{xu2019,schuff2020,xiao2022}
and to estimate several parameters at once in a photonic experiment \cite{cimini2023}. Neural networks trained by an
evolution strategy or reinforcement learning replace the costly expected-gain optimisation of Bayesian design by a fast learned heuristic
\cite{fiderer2021}, and model-aware reinforcement learning trains such a policy through a particle-filter posterior,
with applications to N-V Ramsey magnetometry \cite{belliardo2024quantum,belliardo2024pra}. For a wide field range, a
Bayesian neural network first narrows the range and a reinforcement-learning agent then refines the estimate
\cite{guo2026}.

What is missing from this research is a reference for the ODMR search itself: a bound on how fast any strategy can
locate the resonance. Without it, the reported speedups cannot be placed against the best achievable, it is not
known how many of the measurements are unavoidable, and the observed growth of the speedup as the line occupies less
of the window \cite{dushenko2020sbed} remains unexplained.

Estimation theory provides such references. The Cram\'er--Rao bound limits the variance of an unbiased estimate
by the Fisher information of the data \cite{cramer1946}, but it ignores what is known before the first
measurement. The van Trees inequality, a Bayesian form of the Cram\'er--Rao bound, adds the information of the
prior and holds for any estimator \cite{vantrees1968,gill1995}. A search window is exactly such prior knowledge, so
this bound is a natural yardstick for an ODMR search, and the information it counts is the same quantity that
Bayesian experiment design maximises. In this work we use it to analyse sequential Bayesian experiment design (SBED)
and drive-frequency sweeps for CW and pulsed ODMR.

Section~\ref{sec:theory} derives the bound from the physical model and the speedup, sensitivity and bandwidth that
follow from it; Sec.~\ref{sec:methods} describes the particle filter, the measurement strategies and how they are
evaluated; Sec.~\ref{sec:results} presents the simulations and the experiment; and Sec.~\ref{sec:discussion}
discusses the limits of the approach.

